\documentclass[11pt]{article}
\usepackage[margin=1in]{geometry}
\usepackage{amsmath,amssymb,amsthm,bm,booktabs,hyperref}
\newtheorem{theorem}{Theorem}
\newtheorem{proposition}{Proposition}
\newtheorem{definition}{Definition}
\newtheorem{lemma}{Lemma}
\newtheorem{conjecture}{Conjecture}
\title{CovQ: Compiling Fisher-Information Contracts to Pair-Entangled Quantum
Programs\\\large Hardware Matching Polytopes, Certified Minimum-Cost Schedules, and
Information Floors}
\author{Hina Chauhan Dixit\\Decompute Inc.}
\date{Working draft v0.3}

% Supersedes archive/v0.2/CovQ_Paper_Skeleton_v0.2.tex.
% Basis: docs/proofs/THEOREM_PAIR_WIDTH_AND_INFORMATION_FLOOR_v0.1.md (2A-2F closed),
%        docs/audits/PRIOR_ART_SWEEP_TASK0B5_v0.1.md,
%        docs/decisions/DECISION_001_PIVOT_TO_BOUNDED_WIDTH.md.

\begin{document}
\maketitle

\begin{abstract}
Quantum compilers receive a target unitary, statevector, or circuit. We compile a different
object: a \emph{Fisher-information contract}. The caller supplies a linear map $A$ naming the
parameter combinations it must estimate, a required information floor $G_{\rm req}$, and a
hardware graph $H$; the compiler returns the least-cost quantum program whose realised quantum
Fisher information matrix $F_\Pi$ satisfies $A^{\mathsf T}F_\Pi A\succeq G_{\rm req}$, or a
certificate that no such program exists at the requested entanglement width. Our enabling
result makes the feasible set exact for pair-entangled labelled programs: such a program
realises $F$ on $H$ if and only if $\operatorname{diag}F=\bm1$, $F$ is supported on $E(H)$, and
the matrix of off-diagonal magnitudes lies in the matching polytope of $H$. The
characterisation is constructive---it emits a schedule of parallel Bell pairs, each branch of
two-qubit depth exactly one---and decidable in polynomial time by Edmonds separation, against
NP-hard membership at unbounded width. It turns the contract into a convex program over the
matching polytope whose objective \emph{is} the expected per-shot Bell-pair bill, solved by
cutting planes that return a feasible schedule, a lower bound on achievable cost, or an
infeasibility proof. Two consequences are analytic: a closed-form compilation curve
$\mathrm{cost}^\star(\gamma)=m(\gamma-1)/2$ for a common-mode contract, and an odd/even
ceiling in which a collective mode on $k$ qubits is capped at $2$ for even $k$ but at $2-1/k$
for odd $k$---Edmonds' blossom inequality restated as physics. We show the standard
$k$-producibility witness is the support function of the reachable body in a single direction,
compare against four realisation-matched baselines, and separate labelled schedules from
coherently flagged programs by an explicit dephasing test.
\end{abstract}

\section{Introduction}
\paragraph{Fisher-information contracts, not target states.}
Motivate $(\mathcal G,H,A,G_{\rm req},\mathcal C)\mapsto\Pi$ with
$\min_\Pi\operatorname{Cost}(\Pi)$ subject to $A^{\mathsf T}F_\Pi A\succeq G_{\rm req}$. A user
declares what downstream estimation must be able to resolve; the compiler returns the cheapest
program that certifiably delivers it. State early why exact matching is \emph{not} the
semantics: for unit-diagonal $F,F_\star$, $F\succeq F_\star$ forces $F=F_\star$, so exact
matching, approximate matching and one-sided guarantees collapse into one notion.

\paragraph{Two hazards, cleared before the theorem.}
Matchings are the standard model for \emph{parallel two-qubit gate scheduling}, where a
matching is a time slice inside one circuit; here a matching is a \emph{branch of a randomised
schedule}. And the clique-partitioning polytope is over $0/1$ transitive edge indicators,
whereas our feasible body is over signed correlations in $[-1,1]$ and equals the correlation
polytope at $k=m$.

\paragraph{Novelty boundary.}
Say plainly what is occupied: expectation-value-constrained synthesis (QUEST), correlation
and cut polytopes, matching and size-constrained clique-partitioning polytopes, QFIM
entanglement witnessing, optimal experimental design. Claim only the identification of the
labelled commuting-Pauli program class with a signed \emph{hardware} matching polytope, its
constructive quantum achievability, and the width-two tractable boundary.

\section{Related Work}
\subsection{Expectation-value targeting and partial specification}
\subsection{Correlation, cut, matching, and bounded-cluster polytopes}
\subsection{QFIM and entanglement structure}
\subsection{Optimal experimental design over additive Fisher information}
Frank--Wolfe here is Fedorov--Wynn. Say so.
\subsection{Fisher-information constraints in classical design}
Minimum-cost sensor selection and placement under Fisher-information LMI constraints, and
Loewner-monotone design criteria, already treat information as a \emph{constraint} rather than
an objective. The contract interface is therefore not novel in isolation; what we claim is the
composite with an exact bounded-entanglement achievable set and constructive circuit emission.
\subsection{Matchings in quantum gate scheduling}

\section{Commuting-Pauli Programs, the Fisher IR, and the Resource Vector}
$U_{\bm\theta}=\exp[-\tfrac i2\sum_i\theta_iP_i]$, $F_{ij}=\operatorname{Cov}_\psi(P_i,P_j)$,
$F_{ii}=1-\langle P_i\rangle^2$. Declare the per-shot normalisation once. Define the resource
vector: settings, $k_{\mathrm{logical}}$, $k_{\mathrm{physical}}$, Clifford cost, two-qubit
count and depth, SWAPs, ancillas, and---for coherent programs only---flag dimension and
Schmidt rank.

\section{Background}
\begin{proposition}[Feasibility; inherited]
By symmetric Bernoulli correlation-polytope theory together with the pure-state amplitude
realisation, the feasible zero-mean unit-diagonal QFIMs are exactly
$\operatorname{conv}\{ss^{\mathsf T}\}$.
\end{proposition}
\begin{proposition}[Signed-cat primitive]
$|C_s\rangle=(|s\rangle+|-s\rangle)/\sqrt2$ realises $F=ss^{\mathsf T}$.
\end{proposition}
\begin{lemma}[Labelled additivity]
Per shot, $F_\Pi=\sum_rp_rF_r$. This is a discrete optimal design.
\end{lemma}
\begin{lemma}[Attainability]
$\operatorname{Im}\langle\partial_i\psi|\partial_j\psi\rangle
=-\tfrac14\operatorname{Im}\langle\psi|P_iP_j|\psi\rangle=0$ identically for commuting
Hermitian generators. State the qualifications: local at the operating point; POVM may depend
on it; restrict singular targets to the identifiable support; labelled programs permit
branch-conditioned measurement, coherent programs require joint flag--data measurement;
report an attainable classical Fisher matrix unless the measurement is compiled.
\end{lemma}

\section{Pair-Entangled Compilation and the Hardware Matching Polytope}
\begin{theorem}[Width two, complete graph]
$F\in\mathcal Q^{\mathrm{lab}}_{m,2}$ iff
$(|F_{ij}|)_{i<j}\in\operatorname{MATCH}(K_m)$, i.e.\ iff
$\sum_{j\neq i}|F_{ij}|\le1$ for all $i$ and
$\sum_{\{i,j\}\subseteq S}|F_{ij}|\le\frac{|S|-1}2$ for all odd $S$.
\end{theorem}
\begin{theorem}[Width two, hardware graph]
For $H=(V,E)$ with no routing inside the primitive backend,
\[
\mathcal Q^{\mathrm{lab}}_{H,2}
=\Big\{F:\operatorname{diag}F=\bm1,\;F_{ij}=0\ \text{for}\ ij\notin E,\;
(|F_{ij}|)_{ij\in E}\in\operatorname{MATCH}(H)\Big\}.
\]
Sufficiency is constructive: it emits a schedule of parallel Bell pairs on physical edges.
\end{theorem}
Corollaries: polynomial-time feasibility by Edmonds separation; at most $O(m^2)$ branches by
Carath\'eodory; explicit separating certificates.

\section{Certified Information-Floor Compilation}
\begin{theorem}[Contract compilation]
For non-negative native edge costs, the least-cost pair-entangled labelled program satisfying
$A^{\mathsf T}F_\Pi A\succeq G_{\rm req}$ on $H$ is the value of
\[
\min_{F,t}\ \sum_{e\in E}c_et_e
\quad\text{s.t.}\quad
A^{\mathsf T}FA\succeq G_{\rm req},\
\operatorname{diag}F=\bm1,\
F_{ij}=0\ (ij\notin E),\
|F_{ij}|\le t_{ij},\
t\in\operatorname{MATCH}(H),
\]
and any optimal $F$ is realised by the schedule above at exactly that cost.
\end{theorem}
Down-closedness of $\operatorname{MATCH}(H)$ gives $t_e=|F_e|$ at an optimum, so the objective
is the expected per-shot Bell-pair bill rather than a proxy for it. The only non-polyhedral
constraint separates by an eigenvector, giving an LP-plus-cutting-plane scheme. \textbf{Claim
discipline:} convex, certificate-producing optimisation to prescribed precision---not a
strongly polynomial algorithm.

\paragraph{Certificates.} Exactly one of: a feasible $F$ with its matching decomposition,
emitted circuits and satisfaction margin; a lower bound on achievable cost from a valid
relaxation; or an infeasibility proof for the requested floor.

\section{Operational Consequences}
\begin{proposition}[Odd/even collective-mode ceiling]
For $u_S=\bm1_S/\sqrt k$, width-two compilation caps $u_S^{\mathsf T}Fu_S$ at $2$ when $k$ is
even and at $2-1/k$ when $k$ is odd.
\end{proposition}
Edmonds' blossom inequality as physics: an odd collective mode always leaves one qubit
unmatched, so it provably cannot reach the pair-entanglement ceiling.
\begin{proposition}[Common-mode compilation curve]
For $A=\bm1/\sqrt m$ and $G_{\rm req}=[\gamma]$, at unit edge costs
$\mathrm{cost}^\star(\gamma)=m(\gamma-1)/2$ for $1\le\gamma\le1+2\lfloor m/2\rfloor/m$. A
product probe reaches only $\gamma=1$.
\end{proposition}
\begin{proposition}[Relation to $k$-producibility witnesses]
$\max\{\bm1^{\mathsf T}F\bm1:F\in\mathcal Q^{\rm prog}_{m,k}\}=\lfloor m/k\rfloor k^2+r^2$.
\end{proposition}
The established depth criterion is the support function in one direction; we give the complete
facet description at $k=2$.

\section{Complexity: Width Two versus Width Three}
\begin{conjecture}
Linear utility optimisation over $\mathcal Q^{\mathrm{lab}}_{m,3}$ is NP-hard.
\end{conjecture}
Candidate reduction: Partition into Triangles, utility $1$ on edges and $0$ on non-edges, so
that utility $3q$ on $3q$ vertices is achievable iff the vertices partition into $q$
triangles. \textbf{Unproved, and expected to be a corollary rather than a theorem.} Its role is
to explain why pair width is a privileged tractable regime. Note the structural reason width
three is different: inside a $3$-block the signs obey
$(s_is_j)(s_js_k)(s_is_k)=+1$, so magnitudes alone cannot decide feasibility, and the clean
absolute-value characterisation of Theorem~1 is a genuine feature of the pair level. Compare carefully against the size-constrained
clique-partitioning polytope literature before claiming anything. If this fails and nothing
replaces it, kill gate K6 fires.

\section{Certificates}
Blossom obstructions. The uncertainty-relation reading: for integer $b$ with $\sum_ib_i$ odd,
every realisable $F$ obeys $b^{\mathsf T}Fb=\operatorname{Var}(\sum_ib_iP_i)\ge1$, because
$\sum_ib_iP_i$ has integer spectrum of that parity. A violated inequality is a physical
obstruction, not merely a separating hyperplane. Worked instance: $m=3$, all off-diagonals
$-0.4$; spectrum $\{0.2,1.4,1.4\}$; $b=(1,1,1)$ gives $0.6<1$.

\section{Exact, Sparse, and Approximate Algorithms}
\subsection{Edmonds separation and matching column generation}
\subsection{Full-polytope LP as a background control}
\subsection{Minimum-support formulations and their hardness}
\subsection{Frank--Wolfe with away steps (Fedorov--Wynn) and its duality-gap certificate}

\section{Generator-Aware Clifford Normalisation and Frame-Relative Width}
$k_{\mathrm{logical}}$ versus $k_{\mathrm{physical}}$. A globally entangling Clifford makes
the physical branch entangled; the converse also occurs, so frame choice is a compiler degree
of freedom. Minimum-physical-entanglement language is valid for the physical-$Z$ arm only.

\section{Routing as a Priced Backend}
Routing enlarges the admissible edge set at an explicit SWAP/path cost. Report against the
native-graph feasible set, never folded into it.

\section{Coherent Flags as a Separate Resource Model}
The exact identity $F_\Psi=\sum_rp_rF_r+\operatorname{Cov}_p(\bm\mu_r)$; equality iff branch
means \emph{coincide}. The loophole: $|\Psi_p\rangle=\sum_s\sqrt{p(s)}|s\rangle_{\mathrm
f}|s\rangle_{\mathrm d}$ realises any feasible target with conditional data width one, and
dephasing the flag sends the QFIM to zero. Conclude that conditional branch width is not a
resource here, and give the replacement resource vector. State the comparability rule:
programs on opposite sides of the dephasing test are not comparable on gate count.

\section{Conditioning-Aware Stability}
\begin{proposition}
If $B_\star=A^{\mathsf T}F_\star A\succeq\gamma I$ on its declared support and
$\|A^{\mathsf T}(F-F_\star)A\|\le\varepsilon<\gamma$, then
$\|B^{-1}-B_\star^{-1}\|\le\varepsilon/(\gamma(\gamma-\varepsilon))$. If the kernels differ,
no stable inverse guarantee follows.
\end{proposition}
Report rank, $\lambda^+_{\min}$, $\kappa^+$, kernel principal angles, $\|F^+-F_\star^+\|$.
Note that the polytope's extreme points $ss^{\mathsf T}$ are rank one: feasible extremes,
poor full-field estimation targets.

\section{Experimental Methodology}
\subsection{Benchmark families, including the blossom-triangle and path/star width probes}
\subsection{Four realisation-matched baselines, plus QUEST}
\subsection{Correctness gates C1--C12, with C2 split into C2a/C2b/C2c}
\subsection{Scientific hypotheses, kept separate from the gates}
\subsection{Reproducibility and forward registration}

\section{Results}
Placeholder until Task 0B.5 completes and the pilot thresholds are frozen.

\section{Application Benchmark}
One target derived from Lorentzian light-ray tomography, with the normalisation declared. All
main claims remain application-independent.

\section{Limitations}
Noncommuting generators. Mixed-state QFI. Subunit-diagonal realised probes, which require
joint $(\bm\mu,M)$ moment feasibility rather than $M\in\mathcal Q_m$. Large-scale polytope
complexity. Hardware noise. Compiler optimality. The width-three hardness conjecture.

\section{Conclusion}
Target statement: commuting-Pauli information requirements admit an exact, constructive,
polynomial-time characterisation at pair entanglement on a hardware graph, and that boundary
--- not full-width feasibility --- is where the compiler has something to say.

\end{document}
